% Apresentação de defesa de TCC em Biblioteconomia
% Modelo do ÉguaTex (https://eguatex.github.io), em domínio público pela licença CC0 1.0:
% use, modifique e distribua sem pedir permissão e sem citar a origem.
%
% Compile com pdfLaTeX e BibTeX (o ÉguaTex faz isso sozinho).
% Arquivos: main.tex (os slides) e referencias.bib (bibliografia).
%
% Cada \begin{frame} ... \end{frame} é um slide. Planeje um slide por minuto de fala:
% uma defesa de 20 minutos pede de 15 a 20 slides.

\documentclass[11pt]{beamer}

% ---------- idioma e fontes ----------
\usepackage[T1]{fontenc}
\usepackage[utf8]{inputenc}
\usepackage[brazilian]{babel}
\usepackage{lmodern}

% ---------- tema Madrid com as cores do curso ----------
% Troque os três valores RGB abaixo pelas cores da sua instituição.
\usetheme{Madrid}
\definecolor{petroleo}{RGB}{0,84,98}
\definecolor{areia}{RGB}{214,168,84}
\definecolor{nevoa}{RGB}{233,240,241}
\setbeamercolor{palette primary}{bg=petroleo!85!white, fg=white}
\setbeamercolor{palette secondary}{bg=petroleo!92!black, fg=white}
\setbeamercolor{palette tertiary}{bg=petroleo!70!black, fg=white}
\setbeamercolor{palette quaternary}{bg=petroleo!60!black, fg=white}
\setbeamercolor{structure}{fg=petroleo}
\setbeamercolor{frametitle}{bg=petroleo, fg=white}
\setbeamercolor{title}{bg=petroleo, fg=white}
\setbeamercolor{block title}{bg=petroleo, fg=white}
\setbeamercolor{block body}{bg=nevoa, fg=black}
\setbeamercolor{block title alerted}{bg=areia!80!black, fg=white}
\setbeamercolor{block body alerted}{bg=areia!15, fg=black}
\setbeamercolor{item projected}{bg=areia!80!black, fg=white}
\setbeamertemplate{navigation symbols}{}
\setbeamertemplate{items}[circle]

% ---------- gráficos e referências ----------
\usepackage{booktabs}
\usepackage{pgfplots}
\pgfplotsset{compat=1.18}
% O beamer já carrega o hyperref, por isso o abntex2cite entra direto.
\usepackage[alf]{abntex2cite}

% ---------- sumário no início de cada seção ----------
% A cada \section, um slide mostra o sumário com a seção atual em destaque.
\AtBeginSection[]{
  \begin{frame}{Sumário}
    \tableofcontents[currentsection, hideallsubsections]
  \end{frame}
}

% ---------- instruções de preenchimento ----------
% As notas em cinza explicam o que pôr em cada slide.
% Troque \instrucoestrue por \instrucoesfalse para esconder todas de uma vez.
\newif\ifinstrucoes
\instrucoestrue
\newcommand{\instrucao}[1]{\ifinstrucoes\par\vfill{\color{gray}\scriptsize\itshape #1\par}\fi}

% ---------- dados da capa: preencha aqui ----------
\title[Defesa de TCC]{Estudo de usuários do serviço de referência de uma biblioteca universitária}
\subtitle{Trabalho de Conclusão de Curso}
\author[Nome do autor]{Nome do autor\\[2pt]{\small Orientação: Nome do orientador}}
\institute[UF Exemplo]{Curso de Biblioteconomia\\Universidade Federal Exemplo}
\date[Ano]{Cidade, Ano}

\begin{document}

% ======================================================================
\begin{frame}[plain]
  \titlepage
\end{frame}

% ======================================================================
\section{Introdução}
% ======================================================================
\begin{frame}{Contexto e problema}
  \begin{itemize}
    \item O serviço de referência orienta a busca e o uso da informação.
    \item A procura presencial caiu, mas as dúvidas dos usuários continuam.
    \item \alert{Problema:} que necessidades os usuários levam hoje ao balcão de referência?
  \end{itemize}
  \instrucao{Em 3 ou 4 tópicos, situe o tema e termine com a pergunta de pesquisa. Fale o restante, sem escrever frases longas no slide.}
\end{frame}

\begin{frame}{Objetivos}
  \begin{block}{Objetivo geral}
    Identificar o perfil e as necessidades de informação dos usuários do serviço de referência.
  \end{block}
  \begin{block}{Objetivos específicos}
    \begin{enumerate}
      \item Caracterizar o público atendido.
      \item Classificar as perguntas recebidas.
      \item Propor melhorias no atendimento.
    \end{enumerate}
  \end{block}
  \instrucao{Copie os objetivos do TCC sem mudar os verbos. A banca vai conferir se as conclusões respondem a cada um.}
\end{frame}

% ======================================================================
\section{Referencial teórico}
% ======================================================================
\begin{frame}{As cinco leis da Biblioteconomia}
  \begin{columns}[T]
    \begin{column}{0.55\textwidth}
      \begin{enumerate}
        \item Os livros são para usar.
        \item A cada leitor o seu livro.
        \item A cada livro o seu leitor.
        \item Poupe o tempo do leitor.
        \item A biblioteca é um organismo em crescimento.
      \end{enumerate}
    \end{column}
    \begin{column}{0.4\textwidth}
      \begin{alertblock}{Ligação com o estudo}
        A quarta lei justifica medir o tempo de espera no atendimento \cite{ranganathan1931}.
      \end{alertblock}
    \end{column}
  \end{columns}
  \instrucao{Mostre só os conceitos que sustentam a análise, com o autor de cada um. Um slide por ideia central basta.}
\end{frame}

% ======================================================================
\section{Metodologia}
% ======================================================================
\begin{frame}{Percurso metodológico}
  \begin{table}
    \centering
    \begin{tabular}{@{}ll@{}}
      \toprule
      Etapa & Descrição \\
      \midrule
      Tipo de pesquisa & descritiva, com abordagem quantitativa \\
      Universo & usuários atendidos em um semestre \\
      Instrumento & questionário com 12 perguntas fechadas \\
      Análise & estatística descritiva e categorias de perguntas \\
      \bottomrule
    \end{tabular}
  \end{table}
  \instrucao{Resuma em tabela o tipo de pesquisa, a amostra, o instrumento e a forma de análise. Cite a norma ou o autor do método, se houver.}
\end{frame}

% ======================================================================
\section{Resultados}
% ======================================================================
\begin{frame}{Tipos de pergunta recebidos}
  \begin{columns}[c]
    \begin{column}{0.58\textwidth}
      \begin{tikzpicture}
        \begin{axis}[
          xbar, scale only axis, width=3.4cm, height=4.2cm,
          xmin=0, xmax=50, xlabel={\% das perguntas (dados de exemplo)},
          symbolic y coords={Normalização, Bases de dados, Localização, Empréstimo},
          ytick=data, nodes near coords, nodes near coords align={horizontal},
          tick label style={font=\footnotesize}, label style={font=\footnotesize},
          bar width=9pt, fill=petroleo, draw=none,
          /pgf/number format/use comma,
        ]
          \addplot[fill=petroleo, draw=none] coordinates {(38,Normalização) (27,Bases de dados) (21,Localização) (14,Empréstimo)};
        \end{axis}
      \end{tikzpicture}
    \end{column}
    \begin{column}{0.38\textwidth}
      \begin{itemize}
        \item<2-> Normalização lidera as dúvidas.
        \item<3-> Bases de dados vêm em seguida.
      \end{itemize}
    \end{column}
  \end{columns}
  \instrucao{Um gráfico por slide, com a fonte dos dados. Revele as conclusões com \texttt{<2->} para guiar a fala.}
\end{frame}

% ======================================================================
\section{Considerações finais}
% ======================================================================
\begin{frame}{Considerações finais}
  \begin{itemize}
    \item O objetivo geral foi atingido: o perfil e as necessidades foram descritos.
    \item A maior parte das dúvidas trata de normalização segundo a NBR 6023 \cite{abnt6023}.
    \item \alert{Proposta:} oficinas mensais de normalização para alunos de graduação.
  \end{itemize}
  \instrucao{Responda a cada objetivo em uma linha, diga as limitações e sugira um trabalho futuro.}
\end{frame}

\begin{frame}[allowframebreaks]{Referências}
  \small
  \bibliography{referencias}
\end{frame}

\begin{frame}[plain]
  \centering\vfill
  {\Large\usebeamercolor[fg]{structure} Obrigado!}\par\bigskip
  Nome do autor\par
  {\small email@exemplo.com}
  \vfill
\end{frame}

\end{document}
